\documentclass[11pt, letterpaper]{article}

\topmargin 0in \textheight 9in \textwidth 6.7in \evensidemargin -.21in \oddsidemargin -.21in \columnsep 0.3in

\usepackage{amsmath}
\usepackage{amsfonts}
\usepackage{amssymb}
\usepackage{graphicx}
\usepackage{array}
\usepackage{booktabs}
\usepackage{enumitem}
\usepackage{hyperref}
\usepackage{xcolor}
\usepackage{natbib}
\usepackage{float}
\usepackage{tabularx}

\hypersetup{
    colorlinks=true,
    linkcolor=blue!70!black,
    urlcolor=blue!70!black,
    citecolor=blue!70!black
}

\setlength{\parindent}{0pt}
\setlength{\parskip}{6pt}

\begin{document}

\pagestyle{empty}

~\vskip1em
{\Large
\noindent \textbf{ECE 657D: Research Project} \\[4pt]
}
{\large
\noindent \textbf{Knowledge Sovereignty vs.\ Contextual Interference in RAG Systems} \\[4pt]
\noindent Meihan Liu \quad | \quad 21220972 \\[4pt]
\noindent University of Waterloo \quad | \quad Winter 2026
}

\vskip 1em
\hrule
\vskip 1em

% =============================================================================
\section*{Abstract}
% =============================================================================

Retrieval-Augmented Generation (RAG) systems depend on retrieved context to ground model outputs, but what happens when that context directly contradicts the model's own parametric knowledge? In this work, we study how two open-weight LLMs---Llama-3-8B and Qwen2.5-7B---handle such conflicts across 2,100 samples stratified by entity saliency. We find that behavior varies dramatically with saliency: for well-known entities, both models tend to stick with their internal answers (around 79\% persistence), while for obscure, long-tail entities, persistence drops below 5\% and the models either follow the false context or hedge. A zero-context probe confirms that verified parametric recall is a strong correlate of conflict resistance. Instruction-level ablation on 100 high-confidence samples shows only modest sensitivity to prompt wording, with persistence ranging from 81\% to 88\% across three instruction variants. These results are consistent across both models, suggesting that the observed patterns reflect general properties of mid-scale LLMs rather than architecture-specific quirks.

% =============================================================================
\section*{Introduction}
% =============================================================================

Retrieval-Augmented Generation (RAG) has become a go-to approach for reducing hallucinations in LLMs by supplying relevant documents at inference time~\citep{lewis2020retrieval}. The basic idea is straightforward: instead of relying solely on what the model ``remembers'' from pre-training, we give it external evidence to work with. In most cases, this works well. But it raises an uncomfortable question: what does the model do when the retrieved context is \textit{wrong}?

This is not a hypothetical concern. In real deployments, retrieval pipelines can surface outdated, noisy, or adversarially manipulated documents. When such documents contradict the model's internal knowledge, the model faces a fundamental tension between trusting what it ``knows'' and trusting what it ``sees.'' Prior work has framed this as a conflict between parametric memory and non-parametric evidence~\citep{chen2022rich, longpre2021}, and has shown that model behavior in these situations is neither fully faithful to context nor fully independent of it~\citep{xie2024adaptive}.

In this project, we try to understand this tension more concretely. We construct a testbed of 2,100 factual questions where the retrieved context has been deliberately altered to contain a wrong answer. The key variable we manipulate is \textit{entity saliency}---roughly, how ``well-known'' the entity in question is. Our hypothesis is simple: the model should be better at resisting false context for facts it has seen many times during pre-training (e.g., ``What year was the Eiffel Tower built?'') than for obscure facts it barely encountered (e.g., the classification of a specific protein).

We test this on two models---Llama-3-8B and Qwen2.5-7B---and find strong, consistent evidence for this saliency gradient. We also use a zero-context probe to independently check whether the model can answer each question without any context, which lets us separate ``informed resistance'' (the model knows the answer and sticks with it) from ``uninformed guessing'' (the model happens to output the right answer for other reasons).

% =============================================================================
\section*{Related Work}
% =============================================================================

The problem of knowledge conflicts in RAG systems has received growing attention. \citet{chen2022rich} characterized the tension between internal parametric priors and external evidence, showing that models do not uniformly defer to retrieved context. \citet{longpre2021} studied entity-based knowledge conflicts in QA and found that model behavior depends heavily on whether the model has strong priors about the entity in question. More recently, \citet{xie2024adaptive} categorized LLM behavior under conflict into a spectrum ranging from ``stubborn persistence'' to ``total adherence,'' observing that the same model can exhibit very different strategies depending on the query.

On the retrieval side, \citet{jeong2024adaptiverag} proposed Adaptive-RAG, which dynamically selects retrieval strategies based on query complexity, reducing unnecessary computation while maintaining accuracy. Their work highlights that not all queries benefit equally from retrieval---a point that connects to our observation that high-saliency facts may not need (and can sometimes be harmed by) external context.

Our work differs from prior studies in a few ways. First, we operate at a larger scale (2,100 stratified samples) than most conflict studies. Second, we include a zero-context parametric probe that independently verifies whether the model can recall each entity, enabling us to distinguish between informed and uninformed conflict responses. Third, we compare two architecturally different models (Llama-3 and Qwen2.5) to check whether the observed patterns generalize, rather than being specific to one model family.

% =============================================================================
\section*{Method}
% =============================================================================

\subsection*{Counterfactual Injection}

We start from the SQuAD dataset, which provides question--context--answer triples. For each sample, we take the gold answer $y_\text{gold}$ within the passage and replace it with a counterfactual distractor $y_\text{fake}$, producing a \textit{conflicting context} $C_\text{conflict}$. The model then receives $C_\text{conflict}$ along with the question and must produce an answer.

To make the perturbations realistic, we use a \texttt{smart\_perturb} function that preserves the type of the original answer: years are replaced with other years, numbers with other numbers, and named entities with plausible-sounding alternatives (prefixed with \texttt{alt\_}). This way, the perturbation is semantic rather than obviously anomalous---the conflicting context still reads naturally.

\subsection*{Saliency Stratification}

We divide entities into three tiers based on how likely the model is to have encountered them during pre-training:

\begin{itemize}[nosep]
    \item \textbf{High-Saliency} ($N=800$): Questions containing keywords like ``president,'' ``capital,'' ``century,'' ``war''---topics that appear frequently in web text.
    \item \textbf{Medium-Saliency} ($N=800$): General SQuAD questions without high-salience indicators.
    \item \textbf{Low-Saliency} ($N=500$): Synthetically constructed questions about long-tail entities drawn from a bank of 200+ unique pairs spanning biology, chemistry, geology, astronomy, and other specialized domains (e.g., \textit{Zfyve26 protein}, \textit{Coesite silica polymorph}).
\end{itemize}

We should note that this stratification is a proxy. The keyword heuristic does not perfectly capture pre-training frequency, and the Low tier uses synthetic entities that differ from High/Medium not just in saliency but also in construction method. We discuss the implications of this in the limitations section.

\subsection*{Zero-Context Parametric Probe}

Before running the conflict experiment, we pose each question to the model \textit{without any context} and record its raw answer ($\hat{y}_\text{base}$). If this answer matches $y_\text{gold}$ (via soft token matching with $\geq$50\% normalized overlap), we mark the sample as \texttt{is\_known\_by\_model=True}. This lets us ask: does the model behave differently in conflict when it demonstrably ``knows'' the answer versus when it does not?

This probe is not a perfect measure of internal knowledge---it is sensitive to prompt phrasing and matching criteria, and a model might ``know'' a fact but fail to surface it under our specific prompt. We treat it as a practical proxy for accessible parametric recall rather than an exhaustive test of knowledge possession.

\subsection*{Behavioral Taxonomy}

We classify each model response into one of three categories using soft-token matching:

\begin{table}[H]
\centering
\caption{Taxonomy of model responses under knowledge conflict.}
\label{tab:taxonomy}
\begin{tabular}{lll}
\toprule
\textbf{Category} & \textbf{Matching Rule} & \textbf{Meaning} \\
\midrule
Persistence & $y_\text{out} \approx y_\text{gold}$ & Model outputs its internal answer, ignoring context. \\
Adherence   & $y_\text{out} \approx y_\text{fake}$ & Model follows the (false) retrieved context. \\
Uncertain   & Neither match                        & Model hedges, refuses, or produces unrelated output. \\
\bottomrule
\end{tabular}
\end{table}

\subsection*{Instruction Ablation}

To test whether prompt wording can shift conflict behavior, we run 100 high-confidence (High-Saliency, probe-correct) samples under three instruction variants: \textbf{Neutral} (standard RAG prompt), \textbf{Strict Context} (``You must ONLY use the provided context''), and \textbf{Strict Parametric} (``Prioritize factual accuracy'').

\subsection*{Cross-Model Comparison}

We run the entire pipeline on both Llama-3-8B and Qwen2.5-7B using identical data, prompts, and classification logic to check whether the observed patterns are model-specific or generalizable.

% =============================================================================
\section*{Experiments and Results}
% =============================================================================

All experiments use Ollama with \texttt{temperature=0.0} for deterministic output. The same 2,100 samples are used for both models.

\subsection*{Overall Behavior Distribution (Llama-3-8B)}

Table~\ref{tab:baseline} shows the aggregate distribution across all 2,100 samples for Llama-3-8B. The majority of responses (61.0\%) fall into Persistence, meaning the model tends to favor its parametric memory over conflicting context. Adherence accounts for 15.5\%, and the remaining 23.4\% are classified as Uncertain.

\begin{table}[H]
\centering
\caption{Overall behavior distribution for Llama-3-8B ($N=2{,}100$).}
\label{tab:baseline}
\begin{tabular}{l|ccc}
\toprule
\textbf{Category} & \textbf{Count} & \textbf{\%} & \textbf{95\% CI} \\
\midrule
Persistence     & 1282 & 61.0 & [59.0, 63.1] \\
Uncertain/Other &  492 & 23.4 & [21.6, 25.2] \\
Adherence       &  326 & 15.5 & [14.0, 17.1] \\
\bottomrule
\end{tabular}
\end{table}

\subsection*{Saliency-Stratified Analysis}

Table~\ref{tab:saliency} breaks down behavior by saliency tier. The pattern is stark: High and Medium tiers show roughly 78--80\% persistence and only 10--12\% adherence, while the Low tier flips entirely---persistence drops to 4.6\%, adherence rises to 29.2\%, and the majority (66.2\%) of responses are uncertain.

The ``Base Acc'' column reports the fraction of samples where the zero-context probe returned the correct answer. The ``Resilience'' column measures, among those probe-correct samples, how many maintained the correct answer even when presented with conflicting context. For High-Saliency entities, resilience is 89.5\%; for Low-Saliency, it drops to 13.2\%.

\begin{table}[H]
\centering
\caption{Behavior by saliency tier (Llama-3-8B). Base Acc = zero-context probe accuracy. Resilience = persistence rate among probe-correct samples.}
\label{tab:saliency}
\begin{tabular}{l|cc|ccc|r}
\toprule
\textbf{Tier} & \textbf{Base Acc} & \textbf{Resilience} & \textbf{Persist} & \textbf{Adhere} & \textbf{Uncert} & \textbf{N} \\
\midrule
High   & 20.2\% & 89.5\% & 79.6\% & 10.9\%  & 9.5\%  & 800 \\
Medium & 17.2\% & 86.2\% & 77.8\% & 11.6\% & 10.6\% & 800 \\
Low    & 7.6\%  & 13.2\% & 4.6\%  & 29.2\% & 66.2\% & 500 \\
\bottomrule
\end{tabular}
\end{table}

\subsection*{Instruction Ablation}

Table~\ref{tab:ablation} reports results for three instruction variants tested on 100 High-Saliency samples. Persistence ranges from 81.0\% (neutral) to 88.0\% (strict parametric), a difference of 7 percentage points. Adherence stays between 4--6\% across all variants. The ``strict context'' instruction, which explicitly tells the model to ignore prior knowledge, does not substantially increase adherence---it only reduces uncertainty while persistence remains at 83\%.

These results suggest that, at least for high-confidence facts, prompt-level instruction has limited ability to override the model's tendency to persist. However, we should be cautious about generalizing this finding, since the ablation was conducted on a single tier (High-Saliency) using one model and one prompt template family.

\begin{table}[H]
\centering
\caption{Instruction ablation on 100 High-Saliency samples (Llama-3-8B).}
\label{tab:ablation}
\begin{tabular}{l|ccc}
\toprule
\textbf{Instruction Variant} & \textbf{Persist (\%)} & \textbf{Adhere (\%)} & \textbf{Uncert (\%)} \\
\midrule
Neutral            & 81.0 & 4.0 & 15.0 \\
Strict Context     & 83.0 & 6.0 & 11.0 \\
Strict Parametric  & 88.0 & 4.0 &  8.0 \\
\bottomrule
\end{tabular}
\end{table}

\subsection*{Interaction: Saliency $\times$ Parametric Knowledge}

Table~\ref{tab:interaction} shows behavior conditioned on both saliency and probe outcome. The most informative comparison is within each tier: in the High tier, probe-correct samples show 89.5\% persistence versus 77.1\% for probe-incorrect---a +12.4\% gap attributable to verified recall. In the Low tier, even probe-correct samples ($N=38$) show only 13.2\% persistence and a striking 50.0\% adherence, suggesting that for obscure entities, even when the model ``knows'' the answer, it is easily overridden by context. We note that the Low-Correct cell has a small sample size and these percentages should be interpreted with appropriate caution.

\begin{table}[H]
\centering
\caption{Behavior by saliency and probe outcome (Llama-3-8B).}
\label{tab:interaction}
\begin{tabular}{ll|ccc|r}
\toprule
\textbf{Saliency} & \textbf{Probe} & \textbf{Persist (\%)} & \textbf{Adhere (\%)} & \textbf{Uncert (\%)} & \textbf{N} \\
\midrule
High   & Correct   & 89.5 & 5.6  & 4.9  & 162 \\
High   & Incorrect & 77.1 & 12.2 & 10.7 & 638 \\
\midrule
Medium & Correct   & 86.2 & 9.4  & 4.3  & 138 \\
Medium & Incorrect & 76.0 & 12.1 & 11.9 & 662 \\
\midrule
Low    & Correct   & 13.2 & 50.0 & 36.8 & 38$^*$ \\
Low    & Incorrect & 3.9  & 27.5 & 68.6 & 462 \\
\bottomrule
\multicolumn{6}{l}{\footnotesize $^*$Small sample size; interpret with caution.}
\end{tabular}
\end{table}

\subsection*{Aggregate Knowledge Conditioning}

Table~\ref{tab:conditioning} pools across tiers. Probe-correct samples show 79.6\% persistence versus 57.5\% for probe-incorrect, confirming that verified parametric recall is associated with stronger conflict resistance. Interestingly, even probe-incorrect samples show substantial persistence (57.5\%), which may reflect distributional familiarity from pre-training---the model may have seen these entities enough to develop some implicit prior, even if it cannot produce the correct answer in a zero-context setting.

\begin{table}[H]
\centering
\caption{Behavior conditioned on parametric knowledge, pooled across tiers (Llama-3-8B).}
\label{tab:conditioning}
\begin{tabular}{l|ccc|r}
\toprule
\textbf{Probe Outcome} & \textbf{Persist (\%)} & \textbf{Adhere (\%)} & \textbf{Uncert (\%)} & \textbf{N} \\
\midrule
Correct   & 79.6 & 12.1 & 8.3  & 338 \\
Incorrect & 57.5 & 16.2 & 26.3 & 1762 \\
\bottomrule
\end{tabular}
\end{table}

\subsection*{Cross-Model Comparison}

Table~\ref{tab:crossmodel} compares Llama-3-8B and Qwen2.5-7B side by side. The two models show remarkably similar behavior across all three tiers. In the High tier, both models maintain about 79\% persistence and 10--11\% adherence. In the Low tier, both collapse to below 5\% persistence, though Qwen shows slightly higher uncertainty (75.0\% vs.\ 66.2\%) and lower adherence (22.2\% vs.\ 29.2\%) than Llama. Overall distributions are nearly identical: 61.0\% vs.\ 60.1\% persistence, 15.5\% vs.\ 14.0\% adherence.

This cross-model consistency is notable. Despite architectural differences (Llama uses grouped-query attention; Qwen uses a different tokenizer and training mix), both models exhibit the same saliency-dependent pattern. This suggests that the observed behavior is a general property of mid-scale LLMs rather than an artifact of any particular model's training.

\begin{table}[H]
\centering
\caption{Cross-model comparison: Llama-3-8B vs.\ Qwen2.5-7B ($N=2{,}100$ each).}
\label{tab:crossmodel}
\begin{tabular}{l|ccc|ccc}
\toprule
 & \multicolumn{3}{c|}{\textbf{Llama-3-8B}} & \multicolumn{3}{c}{\textbf{Qwen2.5-7B}} \\
\textbf{Tier} & \textbf{Persist} & \textbf{Adhere} & \textbf{Uncert} & \textbf{Persist} & \textbf{Adhere} & \textbf{Uncert} \\
\midrule
High   & 79.6\% & 10.9\% & 9.5\%  & 78.9\% & 10.5\% & 10.6\% \\
Medium & 77.8\% & 11.6\% & 10.6\% & 77.2\% & 12.4\% & 10.4\% \\
Low    & 4.6\%  & 29.2\% & 66.2\% & 2.8\%  & 22.2\% & 75.0\% \\
\midrule
\textbf{Overall} & 61.0\% & 15.5\% & 23.4\% & 60.1\% & 14.0\% & 25.9\% \\
\bottomrule
\end{tabular}
\end{table}

% =============================================================================
\section*{Discussion}
% =============================================================================

\textbf{Saliency as the dominant variable.} The most consistent finding across both models is the strong association between entity saliency and conflict behavior. High-Saliency entities show roughly 80\% persistence; Low-Saliency entities show less than 5\%. This gradient is stable across two different model architectures and holds regardless of whether we condition on probe outcome. While we cannot claim a clean causal relationship---our saliency tiers differ in construction method as well as estimated pre-training frequency---the effect size is large and robust enough to warrant attention.

\textbf{The role of parametric recall.} The zero-context probe provides a useful lens for interpreting conflict behavior. Within each tier, probe-correct samples consistently show higher persistence than probe-incorrect ones ($+12.4\%$ in High, $+10.2\%$ in Medium). This suggests that accessible parametric recall is associated with stronger conflict resistance. However, even probe-incorrect samples in the High tier maintain 77.1\% persistence, indicating that factors beyond explicit recall---such as distributional familiarity or implicit priors---also contribute to the model's tendency to resist context.

\textbf{Uncertainty as a mixed signal.} About 23\% of responses across both models fall into the Uncertain category. Some of these responses contain hedging language that explicitly acknowledges the conflict (e.g., ``According to the text, X, though this may differ from...''), which could be interpreted as a form of calibrated uncertainty. Other uncertain responses are simply extraction failures or unrelated outputs. Without a more fine-grained analysis of the uncertain category, we cannot definitively say whether this hedging represents a useful self-calibration mechanism, but its consistent presence across both models is worth noting.

\textbf{Instruction sensitivity.} The ablation results show that prompt wording has limited impact on conflict behavior for high-confidence facts. Even with explicit instructions to ``ONLY use the provided context,'' persistence only drops by 2 percentage points relative to ``prioritize factual accuracy.'' This suggests that, at least for well-known facts, the model's internal priors are difficult to override through prompt engineering alone. Whether this extends to other tiers or other prompt families remains an open question.

\textbf{Low-Saliency vulnerability.} The near-total collapse of persistence in the Low tier (4.6\% for Llama, 2.8\% for Qwen) points to a practical vulnerability. In domains where entities are inherently obscure---biomedical research, legal case law, niche historical records---the model has weak or no parametric anchoring and will readily accept whatever the retrieved context says. This makes such domains particularly susceptible to context poisoning, whether through adversarial attacks or simply noisy retrieval.

\textbf{Limitations.} Several caveats apply to these results:

\begin{enumerate}[nosep]
    \item \textbf{Saliency confounds.} Our saliency tiers differ not only in estimated pre-training frequency but also in construction method (SQuAD-derived vs.\ synthetic). The Low tier's behavior may partly reflect domain shift or lexical unfamiliarity rather than saliency alone.
    \item \textbf{Probe limitations.} The zero-context probe uses a specific prompt and matching criterion. A model may ``know'' a fact but fail to surface it under our setup. Base accuracy (7--20\%) is low, meaning most samples are unresolved by the probe; we are effectively conditioning on a small, potentially biased subset.
    \item \textbf{Small cells.} The Low-Saliency probe-correct cell has only $N=38$ samples. Percentage-level conclusions for this cell should be treated as suggestive rather than definitive.
    \item \textbf{Ablation scope.} The instruction ablation covers one tier (High), one model, and three variants. We cannot generalize these results to low-saliency entities or fundamentally different prompt designs.
    \item \textbf{Classification noise.} The soft-match classifier, while reasonable, is imperfect. Some responses labeled ``Uncertain'' may actually contain partial matches to gold or fake answers that fall below the 50\% threshold. A manual annotation study on a random subset would strengthen confidence in the label distribution.
\end{enumerate}

% =============================================================================
\section*{Conclusion}
% =============================================================================

We studied how Llama-3-8B and Qwen2.5-7B handle knowledge conflicts in RAG settings across 2,100 saliency-stratified samples. Three main findings emerge:

\begin{enumerate}[nosep]
    \item \textbf{Saliency shapes conflict behavior.} Both models show strong persistence ($\sim$80\%) for well-known entities and near-total collapse ($<$5\%) for obscure ones. This pattern is consistent across architectures.
    \item \textbf{Parametric recall correlates with resistance.} Samples where the zero-context probe confirms the model can recall the fact independently show meaningfully higher persistence, though the model also exhibits ``uninformed persistence'' that likely reflects distributional priors.
    \item \textbf{Prompt-level intervention has limited reach.} For high-confidence facts, instruction wording shifts persistence by at most 7 percentage points, suggesting that the model's internal priors are not easily overridden by prompt engineering.
\end{enumerate}

These results highlight a practical concern for RAG deployment: in specialized domains with inherently low-saliency entities, the model's lack of parametric anchoring makes it highly susceptible to following incorrect retrieved context. For such domains, additional safeguards---retrieval quality filters, multi-source verification, or uncertainty-aware decoding---may be necessary to ensure reliable outputs.

\section*{Acknowledgments}

This project utilized Claude (Anthropic) for code assistance and experimental workflow optimization. All experimental design, analysis, and interpretation are original work by the author.

\bibliographystyle{plainnat}
\bibliography{references}

\end{document}
